% ============================================================
% Paper 91 (DE): Das Erdős-Gallai-Theorem und offene Probleme
%                in Gradfolgen
% Autor: Michael Fuhrmann
% Specialist Mathematics Project, Build 168
% Datum: 2026-03-12
% ============================================================
\documentclass[12pt,a4paper]{amsart}

\usepackage{amsmath,amssymb,amsthm,mathtools,enumitem,booktabs,array,hyperref}
\usepackage[utf8]{inputenc}
\usepackage[T1]{fontenc}
\usepackage{geometry}
\geometry{margin=2.5cm}

% ---- Theorem-Umgebungen (Deutsch) ----
\newtheorem{theorem}{Theorem}[section]
\newtheorem{lemma}[theorem]{Lemma}
\newtheorem{corollary}[theorem]{Korollar}
\newtheorem{proposition}[theorem]{Proposition}
\newtheorem{satz}[theorem]{Satz}
\newtheorem{korollar}[theorem]{Korollar}
\theoremstyle{definition}
\newtheorem{definition}[theorem]{Definition}
\newtheorem{definition_de}[theorem]{Definition}
\newtheorem{example}[theorem]{Beispiel}
\theoremstyle{remark}
\newtheorem{remark}[theorem]{Bemerkung}
\newtheorem{bemerkung}[theorem]{Bemerkung}
\newtheorem{conjecture}[theorem]{Vermutung}
\newtheorem{vermutung}[theorem]{Vermutung}

% ---- Operatoren ----
\newcommand{\N}{\mathbb{N}}
\newcommand{\Z}{\mathbb{Z}}
\newcommand{\Q}{\mathbb{Q}}
\newcommand{\R}{\mathbb{R}}

% ============================================================
\begin{document}

\title{Das Erd\H{o}s--Gallai-Theorem und offene Probleme\\
in Gradfolgen}

\author{Michael Fuhrmann}
\address{Specialist Mathematics Project, Build 168}
\email{specialist-maths@project.local}
\date{2026-03-12}

\begin{abstract}
Eine endliche Folge nicht-negativer ganzer Zahlen heißt \emph{graphisch}
(oder \emph{realisierbar}), wenn sie die Gradfolge eines einfachen Graphen ist.
Das Erd\H{o}s--Gallai-Theorem (1960) liefert eine vollständige Charakterisierung:
Eine nicht-wachsende Folge $(d_1, \ldots, d_n)$ ist genau dann graphisch, wenn
ihre Summe gerade ist und die Erd\H{o}s--Gallai-Ungleichungen für jedes $k$ gelten.
Dieses elegante Ergebnis --- zunächst von Fulkerson mittels Netzwerkflüssen,
dann durch kombinatorische Argumente (Hakimi) bewiesen --- hat eine reiche Theorie
hervorgebracht.
Diese Arbeit präsentiert den Satz und seine Beweise, den Gradfolgenpolytop,
die gerichtete Variante (Satz von Gale--Ryser), Multigraph-Gradfolgen sowie
offene Probleme: Charakterisierung von Gradfolgen planarer, außerplanarer,
chordaler und Schwellenwertgraphen sowie das Zählproblem für Realisierungen.
\end{abstract}

\maketitle

\tableofcontents

% ============================================================
\section{Einleitung}
% ============================================================

Die \emph{Gradfolge} eines Graphen $G = (V, E)$ mit $V = \{v_1, \ldots, v_n\}$
ist die nicht-wachsend geordnete Folge
\[
  d_1 \geq d_2 \geq \cdots \geq d_n \geq 0,
\]
wobei $d_i = \deg(v_i)$ die Anzahl der inzidenten Kanten von $v_i$ ist.

Eine fundamentale Frage der Graphentheorie ist das inverse Problem:

\begin{definition_de}
Eine Folge $(d_1, \ldots, d_n)$ nicht-negativer ganzer Zahlen heißt
\emph{graphisch} (oder \emph{realisierbar}), wenn ein einfacher Graph
(ohne Schleifen und Mehrfachkanten) auf $n$ Knoten mit dieser Gradfolge existiert.
\end{definition_de}

Die Charakterisierung graphischer Folgen gelang unabhängig Erd\H{o}s--Gallai (1960)
und Hakimi (1962).

\begin{satz}[Erd\H{o}s--Gallai 1960]
\label{satz:EG}
Eine nicht-wachsende Folge nicht-negativer ganzer Zahlen
$d_1 \geq d_2 \geq \cdots \geq d_n$
ist genau dann graphisch, wenn:
\begin{enumerate}[label=(\roman*)]
  \item $\sum_{i=1}^n d_i$ gerade ist, und
  \item Für jedes $k = 1, 2, \ldots, n$:
        \[
          \sum_{i=1}^k d_i \leq k(k-1) + \sum_{i=k+1}^n \min(d_i, k).
        \]
\end{enumerate}
\end{satz}

% ============================================================
\section{Grundbegriffe}
% ============================================================

\begin{definition_de}
Ein \emph{einfacher Graph} ist ein Paar $G = (V, E)$, wobei $V$ eine endliche
Knotenmenge und $E \subseteq \binom{V}{2}$ eine Menge ungeordneter Paare (Kanten)
ohne Schleifen oder Mehrfachkanten ist.
\end{definition_de}

\begin{lemma}[Handschlaglemma]
Für jeden Graphen $G = (V, E)$ gilt:
\[
  \sum_{v \in V} \deg(v) = 2|E|.
\]
\end{lemma}

\begin{proof}
Jede Kante $\{u,v\}$ trägt $1$ zu $\deg(u)$ und $1$ zu $\deg(v)$ bei.
\end{proof}

\begin{korollar}
\label{kor:gerade}
Eine notwendige Bedingung für eine graphische Folge ist, dass ihre Summe gerade ist.
\end{korollar}

% ============================================================
\section{Beweis des Erd\H{o}s--Gallai-Theorems}
% ============================================================

\subsection{Notwendigkeit}

\begin{proposition}[Notwendigkeit der Erd\H{o}s--Gallai-Bedingungen]
Ist $(d_1, \ldots, d_n)$ graphisch, so gelten beide Bedingungen von
Satz~\ref{satz:EG}.
\end{proposition}

\begin{proof}
Bedingung (i) folgt aus dem Handschlaglemma.
Für Bedingung (ii): Sei $G$ eine Realisierung.
Teile $V$ in $S = \{v_1, \ldots, v_k\}$ und $T = V \setminus S$ auf.
Kantenzählung:
\begin{itemize}[nosep]
  \item Kanten innerhalb $S$: höchstens $\binom{k}{2} = k(k-1)/2$, beitragen
        höchstens $k(k-1)$ zu $\sum_{i=1}^k d_i$.
  \item Kanten zwischen $S$ und $T$: Knoten $v_i \in T$ hat höchstens
        $\min(d_i, k)$ Kanten zu Knoten in $S$.
\end{itemize}
Summation liefert $\sum_{i=1}^k d_i \leq k(k-1) + \sum_{i=k+1}^n \min(d_i, k)$.
\end{proof}

\subsection{Hinreichendheit: Hakimis Algorithmus}

\begin{satz}[Hakimi 1962]
\label{satz:hakimi}
Eine nicht-wachsende Folge $(d_1, d_2, \ldots, d_n)$ mit $d_1 \geq 1$ ist
genau dann graphisch, wenn die Folge
\[
  (d_2 - 1, d_3 - 1, \ldots, d_{d_1+1} - 1, d_{d_1+2}, \ldots, d_n)
\]
nach Umsortierung in nicht-wachsender Reihenfolge graphisch ist.
\end{satz}

\begin{proof}
\textbf{Notwendigkeit}: In jeder Realisierung kann man $v_1$ (mit Grad $d_1$)
stets mit $v_2, \ldots, v_{d_1+1}$ verbunden wählen (``Austauschtrick'').
Entfernen von $v_1$ ergibt die reduzierte Folge.

\textbf{Hinreichendheit}: Hat die reduzierte Folge eine Realisierung $G'$,
so füge Knoten $v_1$ hinzu, verbunden mit den $d_1$ Knoten höchsten Grades in $G'$.
\end{proof}

\begin{bemerkung}
Hakimis Satz liefert einen rekursiven (konstruktiven) Beweis der Hinreichendheit.
Zusammen mit der Notwendigkeit beweist dies das Erd\H{o}s--Gallai-Theorem.
Der Algorithmus läuft in $O(n \log n)$ Zeit.
\end{bemerkung}

\begin{example}
Ist $(4, 3, 3, 2, 2)$ graphisch?
Summe $= 14$, gerade. $\checkmark$
EG-Bedingungen für $k=1$: $4 \leq 0 + 1+1+1+1 = 4$. $\checkmark$
$k=2$: $7 \leq 2 + 2+2+2 = 8$. $\checkmark$
$k=3$: $10 \leq 6 + 2+2 = 10$. $\checkmark$
$k=4$: $12 \leq 12 + 2 = 14$. $\checkmark$
Graphisch! Hakimi: Entferne 4, subtrahiere 1 von nächsten 4: $(2,2,1,1)$.
Weiter: $(1,0,1) \to (1,1,0)$; dann $(0,0)$. $\checkmark$
\end{example}

% ============================================================
\section{Fulkersons Bedingungen und Netzwerkflüsse}
% ============================================================

\begin{definition_de}
Die \emph{Fulkerson-Bedingungen} für eine Folge $(d_1, \ldots, d_n)$ lauten:
Für alle disjunkten Teilmengen $S, T \subseteq [n]$:
\[
  \sum_{i \in S} d_i \leq |S|(|S|-1) + 2|S||T|
    + \sum_{j \notin S \cup T} \min(d_j, |S|).
\]
\end{definition_de}

\begin{satz}[Fulkerson 1960]
Die Erd\H{o}s--Gallai-Bedingungen sind äquivalent zu Fulkersons Bedingungen
für den Spezialfall $S = \{1, \ldots, k\}$, $T = \emptyset$.
Ein flusbasierter Beweis liefert die Charakterisierung.
\end{satz}

% ============================================================
\section{Der Gradfolgenpolytop}
% ============================================================

\begin{definition_de}
Der \emph{Gradfolgenpolytop} auf $n$ Knoten ist
\[
  \mathcal{D}_n = \operatorname{conv}\{d(G) : G \text{ einfacher Graph auf } [n]\}
  \subseteq \R^n.
\]
\end{definition_de}

\begin{satz}[Polytopecharakterisierung]
$\mathcal{D}_n$ ist ein Polytop, dessen Ecken genau die Gradfolgen von
Schwellenwertgraphen sind.
Seine Facetten entsprechen den Erd\H{o}s--Gallai-Ungleichungen.
\end{satz}

\begin{bemerkung}
Der Gradfolgenpolytop ist eng mit dem Birkhoff-Polytop und
Transportationspolytopen verwandt.
\end{bemerkung}

% ============================================================
\section{Gerichtete Graphen: Der Satz von Gale--Ryser}
% ============================================================

\begin{definition_de}
Ein Paar von Folgen $(\mathbf{a}, \mathbf{b})$ mit $a_i, b_j \geq 0$ heißt
\emph{bipartit graphisch}, wenn eine $0$-$1$-Matrix $A$ mit Zeilensummen $a_i$
und Spaltensummen $b_j$ existiert.
\end{definition_de}

\begin{satz}[Gale 1957, Ryser 1957]
\label{satz:GR}
Das Paar $\mathbf{a} = (a_1, \ldots, a_m)$ (nicht-wachsend) und
$\mathbf{b} = (b_1, \ldots, b_n)$ ist bipartit graphisch genau dann, wenn
$\sum a_i = \sum b_j$ und für jedes $k = 1, \ldots, m$:
\[
  \sum_{i=1}^k a_i \leq \sum_{j=1}^n \min(b_j, k).
\]
\end{satz}

\begin{bemerkung}
Satz~\ref{satz:GR} ist das bipartite Analogon des Erd\H{o}s--Gallai-Theorems.
Er hat direkte Anwendungen auf das Problem der Realisierung von $(0,1)$-Matrizen
mit vorgegebenen Zeilen- und Spaltensummen.
\end{bemerkung}

% ============================================================
\section{Multigraph-Gradfolgen}
% ============================================================

\begin{definition_de}
Ein \emph{Multigraph} erlaubt Mehrfachkanten zwischen demselben Knotenpaar
(aber keine Schleifen).
Eine Folge $(d_1, \ldots, d_n)$ heißt \emph{multigraphisch}, wenn sie
Gradfolge eines Multigraphen ist.
\end{definition_de}

\begin{satz}[Multigraphische Charakterisierung]
Eine nicht-wachsende Folge $(d_1, \ldots, d_n)$ ist multigraphisch genau dann, wenn:
\begin{enumerate}[label=(\roman*)]
  \item $\sum d_i$ gerade ist;
  \item $d_1 \leq \sum_{i=2}^n d_i$ (der Maximalgrad überschreitet nicht die
        Summe aller übrigen Grade).
\end{enumerate}
\end{satz}

\begin{bemerkung}
Die multigraphische Charakterisierung ist wesentlich einfacher als die
einfach-graphische (Erd\H{o}s--Gallai), da Mehrfachkanten die
$\min(d_i, k)$-Einschränkung aufheben.
\end{bemerkung}

% ============================================================
\section{Schwellenwertgraphen}
% ============================================================

\begin{definition_de}
Ein Graph $G$ heißt \emph{Schwellenwertgraph} (\textit{threshold graph}),
wenn er aus einem einzelnen Knoten aufgebaut werden kann, indem man abwechselnd:
\begin{itemize}[nosep]
  \item einen isolierten Knoten hinzufügt (mit keinem existierenden Knoten verbunden), oder
  \item einen dominierenden Knoten hinzufügt (mit allen existierenden Knoten verbunden).
\end{itemize}
\end{definition_de}

\begin{satz}[Chvátal--Hammer 1977]
Für einen Graphen $G$ sind äquivalent:
\begin{enumerate}[label=(\roman*)]
  \item $G$ ist ein Schwellenwertgraph;
  \item $G$ enthält keinen induzierten Teilgraphen isomorph zu $P_4$, $C_4$ oder $2K_2$;
  \item Die Gradfolge von $G$ ist (bis auf Isomorphie) die einzige Realisierung
        ihrer Gradfolge.
\end{enumerate}
\end{satz}

% ============================================================
\section{Offene Probleme: Planare Graphen}
% ============================================================

\begin{definition_de}
Ein Graph heißt \emph{planar}, wenn er in der Ebene ohne Kantenkreuzungen
gezeichnet werden kann.
\end{definition_de}

\begin{vermutung}[Charakterisierung planarer Gradfolgen]
\label{verm:planar}
Charakterisiere alle Folgen $(d_1, \ldots, d_n)$, die Gradfolge eines
\emph{planaren} Graphen sind.
\end{vermutung}

Vermutung~\ref{verm:planar} ist offen: Keine vollständige Charakterisierung
planarer Gradfolgen analog zu Erd\H{o}s--Gallai ist bekannt.

\begin{satz}[Notwendige Bedingungen für planare Gradfolgen]
Ist $(d_1, \ldots, d_n)$ planar-graphisch, so gilt:
\begin{enumerate}[label=(\roman*)]
  \item $\sum d_i \leq 6n - 12$ (aus Eulers Formel: planarer Graph mit
        $|E| \leq 3n-6$);
  \item Die Folge erfüllt die Erd\H{o}s--Gallai-Bedingungen.
\end{enumerate}
\end{satz}

\begin{bemerkung}
Die Schwierigkeit liegt darin, dass Planarität eine globale Bedingung ist
(kein $K_5$- oder $K_{3,3}$-Minor, nach Kuratowski/Wagner), während
Erd\H{o}s--Gallai rein lokal ist.
\end{bemerkung}

% ============================================================
\section{Offene Probleme: Chordale und außerplanare Graphen}
% ============================================================

\begin{definition_de}
Ein Graph heißt \emph{chordal}, wenn jeder Kreis der Länge $\geq 4$ eine
Sehne besitzt.
Ein Graph heißt \emph{außerplanar}, wenn er eine planare Einbettung mit allen
Knoten auf der äußeren Fläche hat.
\end{definition_de}

\begin{vermutung}[Chordale Gradfolgen]
\label{verm:chordal}
Charakterisiere Gradfolgen chordaler Graphen.
\end{vermutung}

\begin{vermutung}[Außerplanare Gradfolgen]
\label{verm:ausserplanar}
Charakterisiere Gradfolgen außerplanarer Graphen.
\end{vermutung}

\begin{bemerkung}
Chordale Graphen sind vollkommene Graphen und haben eine perfekte
Eliminationsordnung; ihre Struktur ist gut verstanden, aber die
Gradfolgencharakterisierung bleibt offen.
Für außerplanare Graphen gilt $|E| \leq 2n-3$, also $\sum d_i \leq 4n-6$.
\end{bemerkung}

% ============================================================
\section{Zählen von Realisierungen}
% ============================================================

\begin{definition_de}
Für eine graphische Folge $\mathbf{d} = (d_1, \ldots, d_n)$ sei
$R(\mathbf{d})$ die Anzahl nicht-isomorpher einfacher Graphen mit
Gradfolge $\mathbf{d}$.
\end{definition_de}

\begin{vermutung}[Zählen von Realisierungen]
\label{verm:zaehlen}
Bestimme für eine generische graphische Folge $\mathbf{d}$ die Anzahl
$R(\mathbf{d})$ oder eine asymptotische Formel.
\end{vermutung}

\begin{satz}[McKay--Wormald 1991]
Für Gradfolgen mit Maximalgrad $d_{\max} = O(m^{1/4})$ gilt:
\[
  |\{G : d(G) = \mathbf{d}\}| \sim \frac{(\sum d_i)!}{\prod_i d_i! \cdot 2^m}
  \cdot \exp\left(-\frac{\lambda^2}{4}\right),
\]
wobei $m = \sum d_i/2$ und $\lambda = \sum d_i(d_i-1)/m$.
\end{satz}

\begin{bemerkung}
Das exakte Zählen von Realisierungen ist im Allgemeinen $\#P$-vollständig,
also für große Folgen rechnerisch nicht handhabbar.
\end{bemerkung}

% ============================================================
\section{Ganzzahlige Programmierung und Polytope}
% ============================================================

\begin{satz}[IP-Formulierung]
Eine Folge $\mathbf{d}$ ist graphisch genau dann, wenn das folgende
ganzzahlige lineare Programm eine Lösung hat:
\[
  \text{finde } x_{ij} \in \{0,1\} \quad \text{für } i < j \leq n
  \quad \text{mit} \quad \sum_{j \neq i} x_{ij} = d_i \; \forall i.
\]
\end{satz}

\begin{bemerkung}
Das Erd\H{o}s--Gallai-Theorem ist äquivalent dazu, dass die LP-Relaxierung
dieses Programms stets ganzzahlige optimale Lösungen hat, wenn die rechte Seite
die EG-Bedingungen erfüllt.
\end{bemerkung}

% ============================================================
\section{Zusammenfassung und offene Probleme}
% ============================================================

\begin{satz}[Erd\H{o}s--Gallai, Zusammenfassung]
Eine nicht-wachsende Folge $(d_1, \ldots, d_n)$ nicht-negativer ganzer Zahlen ist
die Gradfolge eines einfachen Graphen genau dann, wenn:
\[
  \sum_{i=1}^n d_i \equiv 0 \pmod{2}
  \quad \text{und} \quad
  \sum_{i=1}^k d_i \leq k(k-1) + \sum_{i=k+1}^n \min(d_i, k)
  \quad \forall\, k.
\]
\end{satz}

\noindent\textbf{Offene Probleme:}
\begin{enumerate}[label=(\arabic*)]
  \item \textbf{Planare Gradfolgen}: Kein Analogon zu Erd\H{o}s--Gallai bekannt.
  \item \textbf{Außerplanare Gradfolgen}: Vollständige Charakterisierung offen.
  \item \textbf{Chordale Gradfolgen}: Vollständige Charakterisierung offen.
  \item \textbf{Zählen}: Effizientes exaktes Zählen von $R(\mathbf{d})$.
  \item \textbf{Flächen}: Gradfolgen von Graphen auf Flächen vom Geschlecht $g$.
  \item \textbf{Gleichmäßige Stichproben}: Effizientes Abtasten einer
        gleichmäßigen Verteilung über alle Realisierungen.
\end{enumerate}

% ============================================================
\section{Schluss}
% ============================================================

Das Erd\H{o}s--Gallai-Theorem liefert eine elegante und effiziente Charakterisierung
von Gradfolgen einfacher Graphen.
Der 1960 bewiesene und 1962 durch Hakimis konstruktiven Ansatz verfeinerte Satz
verbindet Netzwerkflusstheorie, Polytopgeometrie und kombinatorische Zählprobleme.

Trotz dieser vollständigen Lösung für einfache Graphen bleiben die analogen Probleme
für strukturell eingeschränkte Graphen (planar, chordal, außerplanar) offen ---
ein eindrucksvolles Beispiel dafür, wie eine lokale kombinatorische Bedingung
(Gradfolge) globale Struktureigenschaften (Planarität) nicht in eine saubere
Charakterisierung fasst.

% ============================================================
\begin{thebibliography}{99}

\bibitem{ErdosGallai1960}
P.\ Erd\H{o}s und T.\ Gallai,
\textit{Graphs with prescribed degrees of vertices} (Ungarisch),
Mat.\ Lapok \textbf{11} (1960), 264--274.

\bibitem{Hakimi1962}
S.\ L.\ Hakimi,
\textit{On realizability of a set of integers as degrees of the vertices of a
linear graph, I},
J.\ SIAM Appl.\ Math.\ \textbf{10} (1962), 496--506.

\bibitem{Fulkerson1965}
D.\ R.\ Fulkerson,
\textit{Zero-one matrices with zero trace},
Pacific J.\ Math.\ \textbf{12} (1965), 831--836.

\bibitem{Gale1957}
D.\ Gale,
\textit{A theorem on flows in networks},
Pacific J.\ Math.\ \textbf{7} (1957), 1073--1082.

\bibitem{Ryser1957}
H.\ J.\ Ryser,
\textit{Combinatorial properties of matrices of zeros and ones},
Canad.\ J.\ Math.\ \textbf{9} (1957), 371--377.

\bibitem{ChvatalHammer1977}
V.\ Chvátal und P.\ L.\ Hammer,
\textit{Aggregation of inequalities in integer programming},
Ann.\ Discrete Math.\ \textbf{1} (1977), 145--162.

\bibitem{McKayWormald1991}
B.\ D.\ McKay und N.\ C.\ Wormald,
\textit{Asymptotic enumeration by degree sequence of graphs with degrees
$o(n^{1/2})$},
Combinatorica \textbf{11} (1991), 369--382.

\bibitem{Diestel2010}
R.\ Diestel,
\textit{Graphentheorie},
4.\ Aufl., Springer, 2010.

\bibitem{BondyMurty2008}
J.\ A.\ Bondy und U.\ S.\ R.\ Murty,
\textit{Graph Theory},
Springer, 2008.

\bibitem{Sierksma1991}
G.\ Sierksma und H.\ Hoogeveen,
\textit{Seven criteria for integer sequences being graphic},
J.\ Graph Theory \textbf{15} (1991), 223--231.

\bibitem{Tripathi2003}
A.\ Tripathi und S.\ Vijay,
\textit{A note on a theorem of Erd\H{o}s and Gallai},
Discrete Math.\ \textbf{265} (2003), 417--420.

\end{thebibliography}

\end{document}
